\subsection{The full special image in weight five}

\begin{statement}{D.51}{Proved}{Task D.1 has dimension three for every prime at least five}
For every prime \(\ell\geq5\), the rational weight-five metabelian
leading image of the full profinite group
\(\operatorname{Out}^{\sharp}(\widehat F_2)\) is exactly
\[
 \mathscr I_{5,\ell}
       =(X+Y)\langle X^2,XY,Y^2\rangle_{\mathbf Q_\ell}[x,y],
 \qquad \dim_{\mathbf Q_\ell}\mathscr I_{5,\ell}=3.
 \tag{D.51.1}
\]
There are actual multiplier-one full profinite special automorphisms
whose three leading images are respectively
\((X+Y)X^2[x,y]\), \((X+Y)XY[x,y]\), and
\((X+Y)Y^2[x,y]\), in the derivation convention
\(D_p(x)=0,\ D_p(y)=[p(X,Y)[x,y],y]\).
This settles the choice left open in D.1 and D.21.
\end{statement}

\begin{proof}
We give the actual source, filtration, and coordinate maps.
Let \(q_\ell:\widehat F(x,y)\twoheadrightarrow P_\ell\)
be the maximal pro-\(\ell\) quotient. Its kernel is characteristic,
so every full automorphism induces an actual automorphism of
\(P_\ell\), and every inner automorphism induces the corresponding
inner automorphism. This defines the map on outer groups.
Write \(\gamma_j\) for the closed ordinary lower central series.
The canonical graded comparison sends
\[
 \gamma_jP_\ell/\gamma_{j+1}P_\ell
          \longrightarrow L_j(\mathbf Z_\ell x\oplus\mathbf Z_\ell y),
 \qquad x\longmapsto x,\quad y\longmapsto y,
 \tag{D.51.2}
\]
and sends each group commutator to the same ordered Lie bracket.
The integral Magnus expansion and its unit leading Lyndon terms
identify this map as an isomorphism of graded lattices.
Its rationalization is the free Lie algebra used in Part D.
The subsequent metabelian map is the actual quotient
\(L_\ell\to L_\ell/L_\ell''\), whose commutator module is
\(\mathbf Q_\ell[X,Y][x,y]\), with
\(X=\operatorname{ad}x\), \(Y=\operatorname{ad}y\).

Filter the full outer group by the inverse images of the
pro-\(\ell\) outer filtration. For a representative acting
trivially through degree \(n\), the classes of
\(\psi(v)v^{-1}\) in
\(\gamma_{n+1}P_\ell/\gamma_{n+2}P_\ell\), for
\(v=x,y\), give its degree-\(n\) derivation.
The commutator identity gives the Leibniz rule on the entire
graded Lie algebra. Changing the representative by an inner
automorphism changes this derivation by the corresponding
inner derivation. These observations give exactly the
filtered outer-to-special-derivation map specified at the
beginning of Part D, followed by \textup{(D.51.2)} and the
displayed metabelian quotient.

Here are literal free-Lie witnesses for the three requested
leading coordinates. For a Lyndon word \(w\) in \(x<y\),
write \(\mathsf L(w)\) for its standard bracket: split at
the longest proper Lyndon suffix and take the bracket of
the two recursively bracketed factors. Define
\[
 \begin{aligned}
 h_1={}&\mathsf L(xxxxy)-\mathsf L(xxxyy)+2\mathsf L(xxyxy),\\
 k_1={}&-\mathsf L(xxxyy)+2\mathsf L(xxyxy),\\[2pt]
 h_2={}&-\mathsf L(xxxyy)-\tfrac12\mathsf L(xxyxy)
             +\mathsf L(xxyyy)+\tfrac32\mathsf L(xyxyy),\\
 k_2={}&\mathsf L(xxyyy)+\tfrac32\mathsf L(xyxyy),\\[2pt]
 h_3={}&\mathsf L(xxyyy)-\mathsf L(xyxyy)-\mathsf L(xyyyy),\\
 k_3={}&-\mathsf L(xyyyy).
 \end{aligned}
 \tag{D.51.3}
\]
Expansion in the free associative algebra gives, for each
\(i\), the exact special equation
\[
                         [h_i,y]=[k_i,x+y].
 \tag{D.51.4}
\]
Thus \(D_i(x)=0,D_i(y)=[h_i,y]\) is a special derivation:
its value at \(-x-y\) is \([k_i,-x-y]\).
All coefficients in \textup{(D.51.3)} belong to
\(\mathbf Z_\ell\) for every \(\ell\geq5\).
Under the canonical metabelian quotient,
\[
 h_1\longmapsto(X+Y)X^2[x,y],\quad
 h_2\longmapsto(X+Y)XY[x,y],\quad
 h_3\longmapsto(X+Y)Y^2[x,y].
 \tag{D.51.5}
\]
For example \(\mathsf L(xxxxy)\) maps to \(X^3[x,y]\),
\(\mathsf L(xxxyy)\) to \(-X^2Y[x,y]\),
\(\mathsf L(xxyyy)\) to \(XY^2[x,y]\), and
\(\mathsf L(xyyyy)\) to \(-Y^3[x,y]\).
The other two displayed Lyndon brackets are brackets
of derived elements and hence map to zero. These are
the exact quotient maps giving \textup{(D.51.5)}.

The integral realization proof in D.44.32 applies to
\(D_i\) in this same \(P_\ell\). Briefly, use its
two prescribed peripheral conjugators through the first
required degree, fixing \(x\). At every later degree
the product error is corrected by the actual map
\((u,v)\mapsto[u,y]+[v,-x-y]\), whose image is the
whole integral free-Lie lattice of that degree by
iterated Jacobi. The correcting conjugators converge
in \(P_\ell\), since every finite \(\ell\)-group
quotient kills a sufficiently deep ordinary lower
central subgroup. The resulting endomorphism has
identity Frattini action and is an automorphism.
Thus there are actual positive peripheral automorphisms
\(\psi_i\) of \(P_\ell\), fixing \(x\), with leading
derivation exactly \(D_i\); no nonzero scaling is
needed for \textup{(D.51.3)}.

Use the exact isomorphism
\(\widehat F/R_{\rm nil}(\widehat F)\simeq\prod_pP_p\)
from D.50.20. Put \(\psi_i\) in its \(\ell\)-component
and the identity in every other component. The actual
local peripheral conjugators, together with identities
in the other components, give three elements of this
product. The conjugator at \(x\) is one.
D.50 supplies a full profinite peripheral automorphism
\(\Phi_i\), fixing \(x\) literally, and the exact square
\[
 \begin{array}{ccc}
 \widehat F&\xrightarrow{\Phi_i}&\widehat F\\
 q_\ell\downarrow&&\downarrow q_\ell\\
 P_\ell&\xrightarrow{\psi_i}&P_\ell.
 \end{array}
 \tag{D.51.6}
\]
Apply the characteristic lower central quotients to
this square, then \textup{(D.51.2)}, and finally the
canonical metabelian Lie quotient. The resulting
leading images are exactly \textup{(D.51.5)}.
This supplies three actual full-source elements with
linearly independent rational leading images.

For the upper bound, the special equation for a
normalized representative at weight five is
\([h,y]=[k,x+y]\).
Writing the two metabelian images as
\(p(X,Y)[x,y]\) and \(q(X,Y)[x,y]\) gives
\(Yp=(X+Y)q\).
Since \(Y\) and \(X+Y\) are relatively prime in the
polynomial ring, \(p\) is divisible by \(X+Y\).
It has degree three, so the entire target is contained
in the three-dimensional space in \textup{(D.51.1)}.
There is no remaining inner ambiguity in this
coordinate: an inner derivation with zero \(x\)-value
would have coefficient \(a\) with \(Xa=0\) in the
torsion-free polynomial commutator module, hence
\(a=0\). This is also the exact comparison in D.21.3.
The lower and upper bounds prove \textup{(D.51.1)}.
\end{proof}

The independent rational certificate
\texttt{de\_special\_source\_witnesses.py}, with its
saved JSON and output, expands every associative
coefficient in \textup{(D.51.3)--(D.51.4)} and applies
the right-\(y\) Fox map
\(M(h)=(h_y y)^{\rm ab}=XYp\) to verify
\textup{(D.51.5)}. The rescalings of the original
source witnesses are exactly \(2,3/2,-1\); their
denominators are respectively \(1,2,1\).
The group realization and transfer are the actual
maps proved above and in D.50.
